\chapter*{\textcolor{mycolor}{Résumé}}
\addcontentsline{toc}{chapter}{Résumé}

Les assistants personnels actuels attendent qu'on les sollicite : il faut poser une question, programmer un
rappel ou ouvrir une application. Or c'est justement au moment où l'on en aurait besoin, en partant de chez soi
par exemple, qu'on n'y pense pas.

Le projet \textbf{LIFEWISE} propose un assistant contextuel qui observe la situation et intervient de lui-même.
Un boîtier équipé d'une caméra et de capteurs (mouvement, porte, luminosité) recueille des indices sur
l'environnement. Un modèle vision-langage léger (VLM) décrit ce que voit la caméra ; un petit modèle de langage
(SLM), précédé de règles écrites, croise cette description avec le contexte (agenda, habitudes) et décide s'il
faut alerter l'utilisateur sur son téléphone. Les événements sont conservés dans une base SQLite locale et
consultables dans un tableau de bord web. Les traitements restent en local, sur un Raspberry Pi qui sert de banc
d'essai pour l'Edge AI.

Le prototype vise un scénario pilote : l'assistance à la sortie du domicile. Au moment du départ, LIFEWISE vérifie
que le téléphone, les clés et le portefeuille sont emportés et, si l'un d'eux n'est pas visible, invite
l'utilisateur à vérifier. L'alerte reste une invitation : une caméra prouve seulement qu'un objet n'est pas dans
son champ, jamais qu'il a été oublié.

Le projet est mené par treize élèves ingénieurs organisés en six sous-équipes, une par maillon de la chaîne.
Ce rapport en présente le contexte, l'état de l'art, la conception, la réalisation et l'évaluation.

\vspace{1cm}

\noindent\textbf{Mots-clés} : vision par ordinateur, Edge AI, Internet des objets, modèle vision-langage,
petit modèle de langage, assistance proactive, Raspberry Pi.
